\begin{abstract}
When carbon-aware schedulers defer GPU jobs to cleaner hours, the deferred jobs that end up emitting
more than if they had run at once are harmed mainly by contention among deferred jobs, not by
carbon-forecast error. Even a clairvoyant greedy scheduler, which knows future carbon intensity exactly
but places jobs one at a time, harms \res{V2U}{ALL-S12-utiltarget0.9}{P5}{harmrate}\%,
\res{V2U}{ALL-S12-utiltarget0.7}{P5}{harmrate}\% and \res{V2U}{ALL-S12-utiltarget0.5}{P5}{harmrate}\% of
the jobs it shifts at 90\%, 70\% and 50\% cluster utilisation. We establish this by replaying the Alibaba
PAI GPU trace against six grids, including the GB system operator's own 48-hour forecasts, and by
attributing each harmed job's excess emissions exactly to forecast error, re-planning and contention.
Contention accounts for \res{V2U}{POOL-S12-utiltarget0.7}{P2}{sharecontention}\% of a point-forecast
scheduler's harmful excess at 70\% utilisation. Conformal, uncertainty-aware policies do not remove
harm; they trade savings for harm along a dial. Their calibration bounds the rate of harmful deferrals
over all decisions, but not the harm rate among the jobs they defer. A congestion-aware rule that keeps
headroom free for later arrivals cuts the harm rate by \res{V2X}{ALL-u0.7}{P8}{harmcutvsPthree}\% for
\res{V2X}{ALL-u0.7}{P8}{savinglossvsPthree} percentage points of saving. On the operator's forecasts it
cuts harm by \res{V2X}{GBop-u0.7}{P8}{harmcutvsPthree}\% at equal savings.
\end{abstract}

\begin{IEEEkeywords}
carbon-aware scheduling, green HPC, GPU clusters, forecast uncertainty, conformal prediction, workload shifting
\end{IEEEkeywords}
